\slajdid{02-s023}
\begin{frame}[label=02-s023]
\gusttekst
Preporučuje se ovakav način crtanja i na ispitu.

Kod crtanja električnih šema tačke kao spojevi se ne preporučuju. Ukrštanje dve linije u tom slučaju se ne smatra spojem.

Ono što treba uočiti
\begin{enumerate}
\item Na ovaj način je moguće realizovati bilo koju kombinacionu mrežu sa bilo kojim brojem ulaza i izlaza.
\item \alert{Principski funkcija realizovana kao zbir proizvoda i realizovana kao proizvod zbirova će zahtevati različit broj resursa.} U ovom primeru funkcija realizovana kao zbir proizvoda zahteva 3 troulazna I kola i jedno troulazno ILI kolo odnosno zahteva 3x3+1x3=12 ulaza. Funkcija realizovana kao proizvod zbirova zahteva 5 troulaznih ILI kola i jedno petoulazno I kolo, odnosno zahteva 5x3+1x5=20 ulaza. Očigledno je u ovom primeru „isplativije“ realizovati funkciju u obliku zbira proizvoda. Principski koja će realizacija biti jednostavnija zavisi od broja logičkih jedinica i nula koje su pojavljuju u funkcionalnoj tabeli, koloni, za željenu funkciju F, kao i kako se pojavljuju.
\item Za realizaciju bilo koje kombinacione mreže su pored invertora dovoljna samo I i ILI logička kola.
\item Bilo koja kombinaciona mreža će, zanemarujući invertore, uvek biti realizovana u maksimalno dva nivoa, odnosno imati „minimalno“ kašnjenje.
\end{enumerate}
\end{frame}
